\section{Datos de preentrenamiento: más allá del volumen, origen, secuencia y memoria externa}

El capítulo anterior explicó cómo calcula el modelo; este pregunta de qué distribución provienen realmente los gradientes. Para la distribución de datos de entrenamiento $q(x)$, el modelo optimiza en la práctica
\[
\mathbb{E}_{x\sim q}\big[\ell_{\theta}(x)\big].
\]
Cambiar $q$ altera los hechos, lenguajes, estilos y patrones de razonamiento que el modelo ve repetidamente. El número de tokens es solo una dimensión; la calidad, la cobertura, la tasa de duplicación, el orden, las interfaces de recuperación y el schedule de crecimiento de parámetros también modifican la supervisión efectiva. La clase divide "los datos son importantes" en cuatro proyectos de investigación comprobables.

\subsection{Los datos no son un barril de petróleo, sino parte del objetivo de entrenamiento}

Llamar a los datos "nuevo petróleo" lleva a pensar solo en la cantidad. Una afirmación más precisa es: los datos definen la distribución de muestreo del riesgo empírico. Si en el corpus web cierto tipo de plantilla, hechos erróneos o proporción lingüística son excesivos, el optimizador los tratará como leyes estables; si las muestras de alto valor son escasas, incluso con suficiente capacidad del modelo, los gradientes no necesariamente reforzarán esa habilidad repetidamente. Por ello, la curaduría de datos, mixture y schedule son componentes externos de la arquitectura.

\lecturefigure{slide-035.jpg}{Impacto de los datos en IA: calidad, volumen, potencia computacional y capacidad acoplados}{CS25 V6 oficial deck, p.~35}

\lecturefigure{slide-036.jpg}{Proyecto de estrategia de datos uno: entrada individual infantil y aprendizaje bilingüe}{CS25 V6 oficial deck, p.~36}

\lecturefigure{slide-037.jpg}{Proyecto de estrategia de datos dos: asignación computacional RAG y escalado guiado por currículo}{CS25 V6 oficial deck, p.~37}

\teachervoice{00:16:20--00:20:23，el docente divide la investigación de datos en cuatro tipos de intervención: cambiar el origen, cambiar la mezcla lingüística, externalizar el conocimiento a recuperación externa y hacer variar conjuntamente la dificultad de los datos con la profundidad del modelo. La diapositiva recuerda: estos proyectos responden preguntas causales distintas y no pueden sintetizarse en una frase "mejores datos, mejor modelo".}

\begin{knowledgebox}{Lectura de la figura: cuatro proyectos que modifican diferentes variables}
Baby Scale modifica la \textbf{fuente de datos familiar/infantil}; Bilingual BabyLM modifica la \textbf{estructura de mezcla lingüística y hablantes}; RAG scaling modifica la \textbf{asignación computacional entre memoria en parámetros y memoria externa}; CGLS modifica el \textbf{orden temporal entre profundidad del modelo y dificultad de las muestras}. Primero se define la intervención, luego se observa el métrica y solo después se pueden discutir conclusiones transferibles.
\end{knowledgebox}

\begin{importantbox}{Los cinco perillas de la estrategia de datos}
\begin{enumerate}
  \item \textbf{Selection}: qué muestras entran al corpus;
  \item \textbf{Mixture}: cómo se proporcionalizan las lenguas, dominios y tareas;
  \item \textbf{Order}: cuándo aparecen los datos simples/difíciles o antiguos/nuevos;
  \item \textbf{Externalization}: si el conocimiento va a los parámetros o al store de recuperación;
  \item \textbf{Co-scaling}: si el schedule de datos varía sincronizado con parámetros, capas y potencia computacional.
\end{enumerate}
Estas perillas modifican los gradientes y no deben verse como una limpieza única antes del entrenamiento.
\end{importantbox}

Estos cuatro proyectos también ilustran el diseño causal básico de la investigación de datos. Si se cambia el origen de los datos, se debe fijar el tokenizer, el número de parámetros, los tokens de entrenamiento y el optimizador; si se cambia la mezcla lingüística, se debe reportar simultáneamente la distribución de prueba por cada idioma; si se cambia el presupuesto de recuperación, se deben incluir el contexto adicional y los costos de índice; si se cambia el schedule de capas, se debe controlar la cantidad total de actualización de parámetros y la potencia computacional de entrenamiento. De lo contrario, el llamado "efecto de datos" podría ser solo que el modelo es más grande, el entrenamiento más largo o la evaluación más cercana a la distribución de entrenamiento. La clase no escribió estos experimentos como una ley universal, sino que proporcionó un método: definir explícitamente un conjunto de intervenciones por vez y escribir los confound sin controlar en los límites de la conclusión.

\subsection{Resumen de la sección}

Los datos de preentrenamiento determinan no solo qué "ve" el modelo, sino también la frecuencia con que se ven las muestras, su orden, contexto y recuperabilidad. Los cuatro casos siguientes responderán por separado: qué aprenden los datos pequeños y naturales, si el bilingüe se interfiere mutuamente, cuánto presupuesto debe tomar RAG y si el crecimiento del modelo puede sincronizarse con el currículo.
